\documentclass[11pt]{article}
\usepackage[margin=1in]{geometry}
\usepackage{amsmath,amssymb}
\usepackage[T1]{fontenc}
\usepackage{lmodern}
\usepackage{hyperref}

\title{Balance-Law Finite-Size SPA Waveform Notes}
\author{}
\date{\today}

\begin{document}
\maketitle

\section{Purpose}

This note documents the finite-size waveform currently implemented in
\texttt{src/finite\_size/waveform.py}.  The important update is that the
stationary-phase approximation (SPA) is not constructed by differentiating a
pre-written closed-form phase.  Instead, the code starts from the binding
energy and flux structure given in Appendix~A of
\href{https://arxiv.org/abs/2410.00294}{arXiv:2410.00294}, solves the
balance-law chirp rate, and then builds both the Fourier phase and Fourier
amplitude from that same chirp rate.

The model is still an inspiral-only benchmark target.  It is not the full
\texttt{IMRPhenomD+FiniteSize} implementation used in the paper.  The goal is a
small, reproducible waveform that tests whether an agent can translate finite
size theory content into a working forecast waveform.

\section{Conventions}

The code uses detector-frame masses.  The chirp mass \(M_c\) is input in solar
masses, the luminosity distance \(d_L\) is input in Mpc, and the total mass is
\begin{equation}
  M = \frac{M_c}{\eta^{3/5}} .
\end{equation}
In seconds,
\begin{equation}
  M_{\rm sec} = M\,M_\odot,
  \qquad
  M_\odot = 4.925491025543576\times 10^{-6}\ {\rm s}.
\end{equation}
The PN velocity variable is
\begin{equation}
  v = (\pi M_{\rm sec} f)^{1/3}.
\end{equation}
The symmetric and antisymmetric spin combinations are
\begin{equation}
  \chi_s = \frac{\chi_1+\chi_2}{2},
  \qquad
  \chi_a = \frac{\chi_1-\chi_2}{2},
  \qquad
  \delta = \sqrt{1-4\eta}.
\end{equation}

The final Fourier-domain waveform convention is
\begin{equation}
  \tilde h(f) =
  A(f)\exp[-i\psi(f)],
  \qquad
  \psi(f) =
  2\pi f t_c-\phi_c-\frac{\pi}{4}+\psi_{\rm bal}(f).
\end{equation}

\section{Energy Balance to SPA}

Appendix~A of arXiv:2410.00294 starts from the balance equation
\begin{equation}
  -{\cal F}_\infty(v)-\dot M(v) = \frac{d{\cal E}}{dt}.
  \label{eq:energy-balance}
\end{equation}
Writing the binding energy as \({\cal E}=M\,e(v)\), with \(M\) constant over
the inspiral benchmark, Eq.~\eqref{eq:energy-balance} gives
\begin{equation}
  \frac{dt}{dv}
  =
  -\,M_{\rm sec}\,
  \frac{de/dv}{{\cal F}_\infty(v)+\dot M(v)} .
  \label{eq:dt-dv}
\end{equation}
For the dominant \((\ell,m)=(2,2)\) mode, \(f=v^3/(\pi M_{\rm sec})\).  Hence
\begin{equation}
  \dot f
  =
  \frac{df}{dt}
  =
  \frac{3v^2}{\pi M_{\rm sec}}\left(\frac{dt}{dv}\right)^{-1}.
  \label{eq:fdot}
\end{equation}
The gravitational-wave phase obeys
\begin{equation}
  \frac{d\Phi_{\rm GW}}{dv}
  =
  2\,\frac{v^3}{M_{\rm sec}}\,\frac{dt}{dv}.
  \label{eq:dphi-dv}
\end{equation}

The code chooses a reference point \(v_{\rm ref}\) at the taper cutoff and
integrates backwards from that point:
\begin{align}
  T(v)
  &=
  \int_v^{v_{\rm ref}} \frac{dt}{dv'}\,dv',
  \\
  \Delta\Phi_{\rm GW}(v)
  &=
  \int_v^{v_{\rm ref}}
  2\,\frac{v'^3}{M_{\rm sec}}\frac{dt}{dv'}\,dv' .
\end{align}
Up to arbitrary constant and linear-in-frequency pieces, which are absorbed by
\(\phi_c\) and \(t_c\), the intrinsic SPA phase is
\begin{equation}
  \psi_{\rm bal}(v)
  =
  \Delta\Phi_{\rm GW}(v)-2\pi f(v)\,T(v).
  \label{eq:balance-phase}
\end{equation}
In the Newtonian limit,
\begin{equation}
  e(v)=-\frac{\eta v^2}{2},
  \qquad
  {\cal F}_\infty(v)=\frac{32}{5}\eta^2 v^{10},
\end{equation}
Eqs.~\eqref{eq:dt-dv}--\eqref{eq:balance-phase} recover the usual curvature
\begin{equation}
  \psi_{\rm bal}(f)
  =
  \frac{3}{128\eta}v^{-5}
  + \hbox{constant}
  + \hbox{linear in } f .
\end{equation}

\section{SPA Amplitude}

The time-domain amplitude is kept at restricted Newtonian order, but the SPA
Jacobian is computed from the same chirp rate as the phase.  The standard
Newtonian TaylorF2 amplitude is
\begin{equation}
  A_N(f)
  =
  \sqrt{\frac{5}{24}}\,
  \frac{M_{c,{\rm sec}}^{5/6}}{d_{L,{\rm sec}}\,\pi^{2/3}}\,
  f^{-7/6}.
\end{equation}
The current balance-law amplitude is
\begin{equation}
  A(f)
  =
  A_N(f)
  \left[
    \frac{\dot f_N}{\dot f}
  \right]^{1/2},
  \qquad
  \dot f_N =
  \frac{96}{5\pi M_{\rm sec}^2}\eta v^{11}.
  \label{eq:balance-amplitude}
\end{equation}
Thus the amplitude reduces exactly to \(A_N\) in the Newtonian limit and
changes only through the finite-size correction to the chirp rate.

\section{Binding Energy Used in the Code}

The paper writes the binding energy as
\begin{equation}
  {\cal E}(v)
  =
  -\frac{M\eta v^2}{2}
  \left[
    {\cal E}_{\rm NS}
    +{\cal E}_{\rm SO}v^3
    +{\cal E}_{\rm SS}v^4
    +{\cal E}_{\rm SSS}v^7
    +{\cal E}_{\rm Love}v^{10}
  \right].
\end{equation}
For the compact benchmark target, the point-particle baseline is kept at
Newtonian order and the spin-induced finite-size pieces are retained:
\begin{equation}
  e(v)
  =
  -\frac{\eta v^2}{2}
  \left[
    1
    +{\cal E}_{\rm SS}^{\rm SIM}(v)v^4
    +{\cal E}_{\rm SSS}^{\rm SIM}v^7
  \right].
  \label{eq:energy-code}
\end{equation}
The quadratic-in-spin piece is
\begin{align}
{\cal E}_{\rm SS}^{\rm SIM}(v)
&=
\chi_a\chi_s
\left(
-\frac{\delta^2\kappa_a}{2}
-\delta\kappa_s
-\frac{\kappa_a}{2}
\right)
\nonumber\\
&\quad
+\chi_s^2
\left(
-\frac{\delta^2\kappa_s}{4}
-\frac{\delta\kappa_a}{2}
-\frac{\kappa_s}{4}
\right)
+\chi_a^2
\left(
-\frac{\delta\kappa_a}{2}
+\eta\kappa_s
-\frac{\kappa_s}{2}
\right)
\nonumber\\
&\quad
+v^2\Bigg[
\chi_a^2
\left(
\frac{25\delta\eta\kappa_a}{12}
-\frac{35\delta\kappa_a}{12}
-\frac{5\eta^2\kappa_s}{6}
+\frac{95\eta\kappa_s}{12}
-\frac{35\kappa_s}{12}
\right)
\nonumber\\
&\qquad
+\chi_a\chi_s
\left(
\frac{25\delta\eta\kappa_s}{6}
-\frac{35\delta\kappa_s}{6}
-\frac{5\eta^2\kappa_a}{3}
+\frac{95\eta\kappa_a}{6}
-\frac{35\kappa_a}{6}
\right)
\nonumber\\
&\qquad
+\chi_s^2
\left(
\frac{25\delta\eta\kappa_a}{12}
-\frac{35\delta\kappa_a}{12}
-\frac{5\eta^2\kappa_s}{6}
+\frac{95\eta\kappa_s}{12}
-\frac{35\kappa_s}{12}
\right)
\Bigg].
\end{align}
The cubic-in-spin piece is
\begin{align}
{\cal E}_{\rm SSS}^{\rm SIM}
&=
\chi_a\chi_s^2
\left(
-2\delta^2\eta\kappa_a
-5\delta^2\kappa_a
+6\delta\eta\kappa_s
-6\delta\kappa_s
-\kappa_a
\right)
\nonumber\\
&\quad
+\chi_s^3
\left(
-\delta^3\kappa_a
-\delta^2\eta\kappa_s
-\frac{9\delta^2\kappa_s}{4}
-\delta\kappa_a
+\frac{\kappa_s}{4}
\right)
\nonumber\\
&\quad
+\chi_a^2\chi_s
\left(
-6\delta\kappa_a
+4\eta^2\kappa_s
+12\eta\kappa_s
-6\kappa_s
\right)
\nonumber\\
&\quad
+\chi_a^3
\left(
-2\delta\eta\kappa_s
-2\delta\kappa_s
+2\eta\kappa_a
-2\kappa_a
\right).
\end{align}

\section{Flux Used in the Code}

The paper writes the flux at infinity as
\begin{equation}
  {\cal F}_\infty
  =
  \frac{32}{5}\eta^2v^{10}
  \left[
    {\cal F}_{\rm NS}
    +{\cal F}_{\rm SO}v^3
    +{\cal F}_{\rm SS}v^4
    +{\cal F}_{\rm SSS}v^7
    +{\cal F}_{\rm Love}v^{10}
  \right].
\end{equation}
The benchmark implementation uses
\begin{equation}
  {\cal F}_\infty
  =
  \frac{32}{5}\eta^2v^{10}
  \left[
    1
    +{\cal F}_{\rm SS}^{\rm SIM}(v)v^4
    +{\cal F}_{\rm SSS}^{\rm SIM}v^7
    +\frac{39}{8}\tilde\Lambda v^{10}
  \right].
  \label{eq:flux-code}
\end{equation}
The \(\tilde\Lambda\) term is an effective leading Love-number flux correction:
with Newtonian binding energy, it reproduces the leading
\(-39\tilde\Lambda v^{10}/2\) tidal phase term quoted in the paper.  A future
production waveform should replace this with the full
\({\cal E}_{\rm Love}\) and \({\cal F}_{\rm Love}\) formulae from
Flanagan--Iyer--Will, as referenced in Appendix~A of the paper.

The quadratic spin-induced flux is
\begin{align}
{\cal F}_{\rm SS}^{\rm SIM}(v)
&=
\chi_a\chi_s(\delta^2\kappa_a+2\delta\kappa_s+\kappa_a)
+\chi_a^2(\delta\kappa_a-2\eta\kappa_s+\kappa_s)
\nonumber\\
&\quad
+\chi_s^2(\delta\kappa_a-2\eta\kappa_s+\kappa_s)
\nonumber\\
&\quad
+v^2\Bigg[
\chi_a^2
\left(
-\frac{127\delta\eta\kappa_a}{16}
+\frac{\delta\kappa_a}{14}
+\frac{43\eta^2\kappa_s}{4}
-\frac{905\eta\kappa_s}{112}
+\frac{\kappa_s}{14}
\right)
\nonumber\\
&\qquad
+\chi_a\chi_s
\left(
-\frac{127\delta\eta\kappa_s}{8}
+\frac{\delta\kappa_s}{7}
+\frac{43\eta^2\kappa_a}{2}
-\frac{905\eta\kappa_a}{56}
+\frac{\kappa_a}{7}
\right)
\nonumber\\
&\qquad
+\chi_s^2
\left(
-\frac{127\delta\eta\kappa_a}{16}
+\frac{\delta\kappa_a}{14}
+\frac{43\eta^2\kappa_s}{4}
-\frac{905\eta\kappa_s}{112}
+\frac{\kappa_s}{14}
\right)
\Bigg]
\nonumber\\
&\quad
+v^3\Bigg[
\chi_a^2(4\pi\delta\kappa_a-8\pi\eta\kappa_s+4\pi\kappa_s)
\nonumber\\
&\qquad
+\chi_a\chi_s(8\pi\delta\kappa_s-16\pi\eta\kappa_a+8\pi\kappa_a)
\nonumber\\
&\qquad
+\chi_s^2(4\pi\delta\kappa_a-8\pi\eta\kappa_s+4\pi\kappa_s)
\Bigg].
\end{align}
The cubic spin-induced flux is
\begin{align}
{\cal F}_{\rm SSS}^{\rm SIM}
&=
\chi_a\chi_s^2
\left(
\frac{13}{3}\delta^2\eta\kappa_a
+\frac{27}{16}\delta^2\kappa_a
+\frac{4}{3}\delta\eta\kappa_s
+\frac{15}{8}\delta\kappa_s
+\frac{3}{16}\kappa_a
\right)
\nonumber\\
&\quad
+\chi_a^2\chi_s
\left(
\frac{95}{12}\delta\eta\kappa_a
+\frac{15}{8}\delta\kappa_a
-\frac{26}{3}\eta^2\kappa_s
+\frac{25}{6}\eta\kappa_s
+\frac{15}{8}\kappa_s
\right)
\nonumber\\
&\quad
+\chi_s^3
\left(
-\frac{7}{4}\delta\eta\kappa_a
+\frac{5}{8}\delta\kappa_a
-\frac{26}{3}\eta^2\kappa_s
-3\eta\kappa_s
+\frac{5}{8}\kappa_s
\right)
\nonumber\\
&\quad
+\chi_a^3
\left(
\frac{29}{6}\delta\eta\kappa_s
+\frac{5}{8}\delta\kappa_s
+\frac{43}{12}\eta\kappa_a
+\frac{5}{8}\kappa_a
\right).
\end{align}

The absorbed flux is
\begin{align}
\dot M(v)
&=
\frac{1}{2}
\left(
9\bar{\cal H}_1^E\eta^2\chi_a
+9{\cal H}_1^E\eta^2\chi_s
\right)v^{15}
\nonumber\\
&\quad
+\frac{1}{2}\Bigg[
\left(
9{\cal H}_1^B\eta^2
+\frac{45}{2}{\cal H}_1^E\eta^2
+\frac{9}{2}\bar{\cal H}_1^E\delta\eta^2
-27{\cal H}_1^E\eta^3
\right)\chi_s
\nonumber\\
&\qquad
+\left(
9\bar{\cal H}_1^B\eta^2
+\frac{45}{2}\bar{\cal H}_1^E\eta^2
+\frac{9}{2}{\cal H}_1^E\delta\eta^2
-27\bar{\cal H}_1^E\eta^3
\right)\chi_a
\Bigg]v^{17}
\nonumber\\
&\quad
+18{\cal H}_0\eta^2v^{18}.
  \label{eq:mdot-code}
\end{align}
Appendix~A prints \(9{\cal H}_0\eta^2v^{18}\).  The code default doubles that
term to implement the BBH electric--magnetic duality convention used by the
benchmark.  Setting \texttt{h0\_dual=False} restores the electric-only
coefficient.

\section{Frequency Band and Taper}

The waveform is inspiral-only.  The cutoff is
\begin{equation}
  f_{\rm ISCO}=\frac{1}{6^{3/2}\pi M_{\rm sec}},
  \qquad
  f_{\rm cut}=q_{\rm max}f_{\rm ISCO},
\end{equation}
where \(q_{\rm max}=\texttt{fmax\_over\_fisco}\), with default \(q_{\rm max}=1\).
The waveform is exactly zero for \(f<f_{\rm low}\) or \(f\ge f_{\rm cut}\).
Inside the band the taper is
\begin{equation}
  W(f)
  =
  \frac{1}{1+\exp[(f-f_{\rm cut})/\sigma]},
  \qquad
  \sigma=
  \texttt{sigma\_taper\_over\_fisco}\,f_{\rm ISCO}.
\end{equation}
The default is \(\texttt{sigma\_taper\_over\_fisco}=0.01\).

\section{Code API}

The benchmark-facing function is
\begin{verbatim}
h_of_f(
    f, Mc, eta, dL,
    chi1=0.0, chi2=0.0,
    tc=0.0, phic=0.0,
    kappa_s=1.0, kappa_a=0.0,
    H0=0.0, Lambda_tilde=0.0,
    H1E=0.0, H1B=0.0,
    H1E_bar=0.0, H1B_bar=0.0,
    f_low=20.0,
    fmax_over_fisco=1.0,
    sigma_taper_over_fisco=0.01,
    phase_only=False,
    h0_dual=True,
)
\end{verbatim}
It returns a complex NumPy array with the same shape as the input frequency
array.  If \texttt{phase\_only=True}, the code returns
\(W(f)\exp[-i\psi(f)]\) inside the band and zero outside it.

\section{Current Caveats}

\begin{itemize}
  \item The point-particle baseline in the balance-law dynamics is Newtonian.
  Higher-PN point-particle terms can be added later by replacing the
  \({\cal E}_{\rm NS}\), \({\cal E}_{\rm SO}\), \({\cal F}_{\rm NS}\), and
  \({\cal F}_{\rm SO}\) blocks.
  \item The Love-number sector is represented by the leading effective flux
  term in Eq.~\eqref{eq:flux-code}.  This keeps the benchmark compact, but it
  is not the full Flanagan--Iyer--Will tidal energy and flux model.
  \item The cutoff is \(f_{\rm ISCO}\)-based, not the paper's
  \(f_{22}^{\rm tape}=0.35f_{22}^{\rm peak}\) prescription.
  \item The implementation is intended for benchmark reproducibility and
  forecasting experiments, not final astrophysical inference.
\end{itemize}

\end{document}
